\chapter{State Space Models}
\label{ch:state-space-models}

\section{What Is A State Space Model?}

A state space model describes a system by separating its internal state from its observed
input and output. The state is not just another feature vector. It is the model's memory
of the past, updated through a dynamical rule.

In continuous time, the standard linear state space model is
\begin{equation}
  \frac{d}{dt}\vx(t) = \mA \vx(t) + \mB \vu(t),
  \qquad
  \vy(t) = \mC \vx(t) + \mD \vu(t).
  \label{eq:continuous-ssm}
\end{equation}
Here $\vu(t)$ is the input signal, $\vx(t)$ is the latent state, and $\vy(t)$ is the
output signal. The matrices have the same conceptual roles as in the discrete recurrence
from Chapter~\ref{ch:linear-systems-and-filters}:
\[
  \mA: \text{state dynamics},\quad
  \mB: \text{input write-in},\quad
  \mC: \text{state read-out},\quad
  \mD: \text{direct input-output path}.
\]

In control theory, this notation is used to model physical systems. In deep learning, we
use the same notation to build sequence layers. The state does not have to correspond to
literal physical variables. It is a learned memory representation.

\section{Why Continuous Time?}

Neural networks usually operate on discrete sequences:
\[
  \vu_1,\ldots,\vu_L.
\]
So it is fair to ask why SSM papers start from continuous-time equations. There are three
main reasons.

\paragraph{Sampling rate should be a parameter.}
Many signals come from continuous processes sampled at some interval. Neural recordings,
audio, sensor streams, and trajectories all have a physical sampling interval. A
continuous-time model can be discretized at different step sizes $\Delta$.

\paragraph{Stability is cleaner in continuous time.}
For the continuous system
\[
  \frac{d}{dt}\vx(t)=\mA\vx(t),
\]
stability is controlled by the real parts of eigenvalues of $\mA$. Negative real part
means decay. This is often easier to parameterize before converting to discrete time.

\paragraph{S4 and HiPPO are naturally continuous.}
The HiPPO theory behind S4 begins as an online continuous-time memory problem. The
continuous formulation is not just decoration; it is part of where the S4 initialization
comes from.

After discretization, however, the model used in a neural network is an ordinary discrete
sequence layer.

\section{Solving The Continuous Linear System}

To discretize \eqref{eq:continuous-ssm}, first understand the solution over one time
interval. Assume the input $\vu(t)$ is held constant over the interval $[t, t+\Delta)$.
This assumption is called zero-order hold. It matches the common situation where a
sampled input value represents the signal value for one bin.

The solution of
\[
  \frac{d}{dt}\vx(t)=\mA\vx(t)
\]
over a time step $\Delta$ is
\[
  \vx(t+\Delta)=e^{\Delta \mA}\vx(t),
\]
where $e^{\Delta \mA}$ is the matrix exponential. This is the matrix analogue of the
scalar solution
\[
  x(t+\Delta)=e^{\Delta a}x(t).
\]

With a constant input over the interval, the solution becomes
\[
  \vx(t+\Delta)
  =
  e^{\Delta \mA}\vx(t)
  +
  \left(\int_0^\Delta e^{\tau \mA}\,d\tau\right)\mB\vu(t).
\]
This formula gives the exact discrete update under zero-order hold.

\section{Discretization}

The discretized SSM has the form
\begin{equation}
  \vx_t = \bar{\mA}\vx_{t-1} + \bar{\mB}\vu_t,
  \qquad
  \vy_t = \mC\vx_t + \mD\vu_t.
  \label{eq:discretized-ssm}
\end{equation}
The barred matrices are discrete-time versions of the continuous parameters. Under
zero-order hold,
\begin{equation}
  \bar{\mA} = e^{\Delta \mA},
  \qquad
  \bar{\mB}
  =
  \left(\int_0^\Delta e^{\tau \mA}\,d\tau\right)\mB.
  \label{eq:zoh-discretization}
\end{equation}
If $\mA$ is invertible, the second expression can be written as
\[
  \bar{\mB}
  =
  \mA^{-1}(e^{\Delta \mA}-\mI)\mB.
\]
The integral form is more general and conceptually safer.

The output matrices $\mC$ and $\mD$ are often left unchanged by discretization in the
standard setup. The dynamics and input write-in need conversion because they describe
how the state evolves over a time interval.

\section{Step Size}

The step size $\Delta$ determines how much continuous time passes between two sequence
positions. If neural data are binned every $20$ ms, then a physically literal choice is
\[
  \Delta = 0.02 \text{ seconds}.
\]
In neural network SSM layers, $\Delta$ is often learned or parameterized per channel or
per state dimension. Then it becomes more than a sampling constant: it controls the
timescale of the learned dynamics.

For a scalar continuous system
\[
  \frac{d}{dt}x(t)=a x(t),
\]
discretization gives
\[
  \bar{a}=e^{\Delta a}.
\]
If $a<0$, then $0<\bar{a}<1$. Larger $\Delta$ makes $\bar{a}$ smaller, meaning faster
decay per discrete step. Smaller $\Delta$ keeps $\bar{a}$ closer to $1$, meaning slower
decay per step.

This is why learned step sizes are powerful. They let different channels or modes operate
at different effective temporal scales.

\section{Continuous And Discrete Stability}

For the continuous autonomous system
\[
  \frac{d}{dt}\vx(t)=\mA\vx(t),
\]
a mode with eigenvalue $\lambda$ evolves like
\[
  e^{\lambda t}.
\]
If $\RePart(\lambda)<0$, it decays. If $\RePart(\lambda)>0$, it grows. If
$\RePart(\lambda)=0$, it persists in magnitude and may oscillate.

After discretization,
\[
  \bar{\lambda}=e^{\Delta \lambda}.
\]
The magnitude is
\[
  |\bar{\lambda}| = |e^{\Delta \lambda}| = e^{\Delta \RePart(\lambda)}.
\]
So
\[
  \RePart(\lambda)<0
  \quad \Longleftrightarrow \quad
  |\bar{\lambda}|<1.
\]
This exactly connects the continuous-time stability condition to the discrete-time unit
circle condition from Chapter~\ref{ch:linear-systems-and-filters}.

\section{Recurrent Mode}

Once discretized, the SSM can be run as a recurrence:
\[
  \vx_t = \bar{\mA}\vx_{t-1} + \bar{\mB}\vu_t,
  \qquad
  \vy_t = \mC\vx_t + \mD\vu_t.
\]
This is the recurrent mode. It is natural for online inference because we only need to
store the current state $\vx_{t-1}$ and the new input $\vu_t$.

For a causal neural decoder, recurrent mode is conceptually attractive:
\[
  \text{new neural bin} \rightarrow \text{update state} \rightarrow \text{emit current representation}.
\]
There is no need to revisit the whole history.

The downside is training parallelism. A naive recurrence computes step $t$ after step
$t-1$, just like a vanilla RNN. SSMs become especially useful when we exploit additional
structure to train them in parallel.

\section{Convolution Mode}

If $\bar{\mA},\bar{\mB},\mC,\mD$ are fixed across time, the discretized SSM is linear and
time invariant. As in Chapter~\ref{ch:linear-systems-and-filters}, unrolling gives
\[
  \vy_t
  =
  \mD\vu_t
  +
  \sum_{i=0}^{t-1}
  \mC\bar{\mA}^i\bar{\mB}\vu_{t-i}.
\]
Define the SSM convolution kernel
\begin{equation}
  \mK_i = \mC\bar{\mA}^i\bar{\mB}.
  \label{eq:discrete-ssm-kernel}
\end{equation}
Then the layer can be evaluated as a causal convolution:
\[
  \vy_t = \mD\vu_t + \sum_{i=0}^{t-1}\mK_i\vu_{t-i}.
\]

This is convolution mode. It exposes parallelism over sequence positions. The layer can
be trained like a convolutional sequence model, even though it also has a recurrent
interpretation.

The central computational challenge is producing or applying the kernel efficiently. A
generic dense $\bar{\mA}$ makes $\bar{\mA}^i$ expensive for long sequences and large state
sizes. S4 and S5 are largely about choosing structures for $\bar{\mA}$ that make this
tractable.

\section{The Role Of The Direct Term}

The matrix $\mD$ gives a direct path from current input to current output:
\[
  \vy_t = \mC\vx_t + \mD\vu_t.
\]
In convolution language, this is an instantaneous kernel component. Depending on indexing
convention, it may be combined with $\mK_0$ or treated separately.

The direct term is useful because the state path may be best at modeling temporal context,
while the direct term preserves immediate information. Many deep architectures also have
residual connections around the whole SSM block, giving another direct route for current
features.

\section{SISO And MIMO}

A single-input single-output SSM has
\[
  \vu_t \in \R,\qquad \vy_t \in \R.
\]
Then $\mB$ is a column vector, $\mC$ is a row vector, and $\mK_i$ is scalar.

A multi-input multi-output SSM has
\[
  \vu_t \in \R^{C_{\mathrm{in}}},
  \qquad
  \vy_t \in \R^{C_{\mathrm{out}}},
\]
and $\mK_i$ is a matrix:
\[
  \mK_i \in \R^{C_{\mathrm{out}} \times C_{\mathrm{in}}}.
\]
This is the natural form for modern sequence models, where every time step has many
features.

S4 often uses SISO SSMs in parallel across channels, followed by separate channel-mixing
layers. S5 emphasizes MIMO SSMs more directly. This is one of the practical simplifications
that makes S5 easier to reason about as a neural network layer.

\section{State Size, Model Width, And Channels}

SSM notation can be confusing because there are several dimensions:
\[
\begin{array}{ll}
L & \text{sequence length}, \\
C_{\mathrm{in}} & \text{input channel count}, \\
C_{\mathrm{out}} & \text{output channel count}, \\
H & \text{model width in a deep network}, \\
N & \text{SSM state size}.
\end{array}
\]
The state size $N$ controls the dimension of the latent dynamical memory. The model width
$H$ is the feature width passed between blocks. In many architectures,
$C_{\mathrm{in}}=C_{\mathrm{out}}=H$ inside the backbone, but $N$ may be different.

This distinction matters for capacity. Increasing $H$ gives the network more feature
channels. Increasing $N$ gives each SSM more dynamical memory. The two changes are not
equivalent.

\section{A Deep SSM Layer Is Not Just The SSM}

The bare SSM recurrence is linear. A practical deep SSM block usually wraps it with other
operations:
\[
  \vz
  \rightarrow
  \text{normalization}
  \rightarrow
  \text{input projection}
  \rightarrow
  \text{SSM temporal mixing}
  \rightarrow
  \text{nonlinearity or gate}
  \rightarrow
  \text{output projection}
  \rightarrow
  \text{residual addition}.
\]
Different papers choose different wrappers, but this pattern is the right mental model.
The SSM supplies long-range temporal mixing. The surrounding neural network supplies
feature mixing, nonlinearity, and optimization stability.

This also explains why a linear SSM can be useful in a nonlinear model. A transformer
attention operation is not the whole transformer block either. It is one structured
mixing operation inside a larger residual architecture.

\section{Why Structure Is Needed}

The formula
\[
  \mK_i = \mC\bar{\mA}^i\bar{\mB}
\]
looks simple, but a naive implementation has problems.

\paragraph{Computing long kernels can be expensive.}
If $\bar{\mA}$ is dense, repeatedly multiplying by $\bar{\mA}$ across a long sequence can
be costly.

\paragraph{Training dense dynamics can be unstable.}
The eigenvalues of $\bar{\mA}$ control memory. Bad parameterization can lead to exploding,
vanishing, or poorly conditioned dynamics.

\paragraph{Long memory needs good initialization.}
Random matrices do not automatically produce a useful spectrum of timescales. For long
sequences, the model needs modes that cover both short and long temporal dependencies.

\paragraph{Parallelism requires exploitable algebra.}
The convolution view is only computationally valuable if the kernel can be generated or
applied efficiently. S4 uses special matrix structure. S5 simplifies toward diagonal
structure and parallel scan. Mamba uses selective scan rather than fixed convolution.

\section{SSMs Versus RNNs}

The recurrence
\[
  \vx_t = \bar{\mA}\vx_{t-1}+\bar{\mB}\vu_t
\]
resembles a vanilla RNN, but there are crucial differences.

A standard RNN has nonlinear state updates:
\[
  \vh_t = \tanh(\mW_h\vh_{t-1}+\mW_u\vu_t).
\]
This nonlinearity makes the recurrence expressive, but it also prevents a simple
convolutional representation.

A linear time-invariant SSM has a linear state update. That sounds less expressive, but
it allows the recurrent and convolutional views to coexist:
\[
  \text{linear recurrence}
  \quad \Longleftrightarrow \quad
  \text{causal convolution}.
\]
The nonlinearity is moved outside the temporal recurrence and placed in the surrounding
deep block. This design is one of the reasons SSM layers can train efficiently on long
sequences.

\section{SSMs Versus CNNs}

An ordinary causal CNN learns kernel slices directly:
\[
  \mK_0,\ldots,\mK_{K-1}.
\]
An SSM generates them:
\[
  \mK_i = \mC\bar{\mA}^i\bar{\mB}.
\]
This gives the SSM a compact way to represent long filters. The price is that the kernel
must have the structure induced by the state dynamics.

So the comparison is:
\[
\begin{array}{lll}
\text{CNN} & \text{free finite kernel} & \text{simple, local unless very large}, \\
\text{SSM} & \text{structured long kernel} & \text{compact, dynamical, stability-sensitive}.
\end{array}
\]
This is the practical reason to study SSMs after CNNs.

\section{SSMs Versus Transformers}

Transformers use content-dependent pairwise interactions:
\[
  \text{attention weight from } t \text{ to } s
  =
  \softmax(q_t^\top k_s).
\]
This gives direct access from one position to another, and the access pattern changes
with the input content. That flexibility is powerful.

Linear time-invariant SSMs do not have this kind of content-dependent routing. The same
kernel is applied regardless of the input values. Their advantage is different: they can
model long temporal filters with linear scaling and recurrent online operation.

Mamba can be understood as an attempt to recover some of the content-dependent behavior
of attention while keeping the linear-time recurrent structure of SSMs. It does this by
making some SSM parameters input-dependent.

\section{What To Remember}

\begin{itemize}
  \item A state space model separates input, latent state, and output.
  \item Continuous-time SSMs are discretized before being used on sampled sequences.
  \item Zero-order hold gives
  $\bar{\mA}=e^{\Delta\mA}$ and
  $\bar{\mB}=(\int_0^\Delta e^{\tau\mA}d\tau)\mB$.
  \item Continuous stability uses negative real eigenvalue parts; discrete stability uses
  eigenvalues inside the unit circle.
  \item A discretized linear time-invariant SSM can run in recurrent mode or convolution
  mode.
  \item The convolution kernel is $\mK_i=\mC\bar{\mA}^i\bar{\mB}$.
  \item S4 and S5 are mostly about making this kernel expressive, stable, and efficient.
\end{itemize}

\section{Exercises}

\begin{enumerate}
  \item For the scalar continuous system $\dot{x}(t)=-2x(t)$, compute the discrete
  multiplier $\bar{a}$ when $\Delta=0.1$.
  \item If a continuous eigenvalue is $\lambda=-1+3i$, what is the magnitude of the
  corresponding discrete eigenvalue after step size $\Delta$?
  \item Explain why zero-order hold is a reasonable assumption for binned neural data.
  \item In your own words, distinguish state size $N$ from model width $H$.
  \item Starting from \eqref{eq:discretized-ssm}, unroll the recurrence for three steps
  and identify $\mK_0,\mK_1,\mK_2$.
\end{enumerate}

\section{Prepares You For}

The next chapter explains why modern SSMs needed special initializations and structured
state matrices to store long histories without becoming unstable or computationally
expensive. That brings us to HiPPO: a principled way to initialize recurrent memory for
long sequence modeling.
